\documentclass[aspectratio=169]{beamer}

\usepackage{hyperref}
\usepackage{qrcode}
\hypersetup{colorlinks,linkcolor=,urlcolor=blue}

%% Theme...
\usefonttheme{serif}

%% Footer...
\defbeamertemplate*{footline}{shadow theme}{
    \leavevmode
    \hbox{
        \begin{beamercolorbox}[
                wd =        0.33\paperwidth,
                ht =        2.5ex,
                dp =        1.125ex,
                leftskip =  0.3cm plus1fil,
                rightskip = 0.3cm
            ]{author in head/foot}
            \flushleft DBS
        \end{beamercolorbox}
        \begin{beamercolorbox}[
                wd =        0.33\paperwidth,
                ht =        2.5ex,
                dp =        1.125ex,
                leftskip =  0.3cm plus1fil,
                rightskip = 0.3cm
            ]{author in head/foot}
            \insertshorttitle
        \end{beamercolorbox}
        \begin{beamercolorbox}[
                wd =        0.33\paperwidth,
                ht =        2.5ex,
                dp =        1.125ex,
                leftskip =  0.3cm plus1fil,
                rightskip = 0.3cm
            ]{title in head/foot}
            \hfill \insertframenumber\,/\,\inserttotalframenumber%
        \end{beamercolorbox}
    }
}

\title{Bases de Datos}
\author{Andrés Oswaldo Calderón Romero, Ph.D.}
\date{\today}

\begin{document}

\frame{\titlepage}

\section{Contenido}

\begin{frame}{Información}
\begin{itemize}
    \item Profesor: Andrés Oswaldo Calderón Romero, Ph.D.
    \item \textbf{Correo}: \href{mailto:andrescalderonr@javeriana.edu.co}{andrescalderonr@javeriana.edu.co} (Iniciar el asunto con ``[DBS]'').
    \item \textbf{Web}: \href{https://www.cs.ucr.edu/~acald013/}{https://www.cs.ucr.edu/$\sim$acald013/}.
    \item Páginas importantes: \url{https://github.com/aocalderon/PUJ/} $\longrightarrow$ \texttt{2026-30}.
    \item Plataformas: 
    \begin{itemize}
        \item BrightSpace (Anuncios y entregas).
        \item GitHub (Material de clase).
    \end{itemize}
    \item Herramientas:  \href{https://www.postgresql.org/}{PostgreSQL}, \href{https://www.pgadmin.org/}{PgAdmin}.
    \item Horario de atención (Office Hours):
    \begin{itemize}
        \item Martes: 10:00 a 11:00am (FIFO).
    \end{itemize}
\end{itemize}
\end{frame}

\section{Los Estudiantes}

\begin{frame}{Los Estudiantes}
\begin{itemize}
    \item ¿Nombre? 
    \item ¿Experiencia en programación? ¿Lenguajes? 
    \item ¿Experiencia en bases de datos? ¿DBMSs?
    \item ¿Algo puntual que espera aprender en este curso?
\end{itemize}
\end{frame}

\section{Sobre el Curso}

\begin{frame}{Sobre el Curso}
\begin{itemize}
    \item ¿Qué sabemos?
    \item ¿Qué nos han dicho?
    \item ¿Qué expectativas tengo?
\end{itemize}

    \centering
    \vspace{10mm}
    {\small \href{https://www.mentimeter.com/app/presentation/aljzuorbyosmeh7v9vvh4ik9e9y3zdcn/edit}{Veámos:}}\\
    \vspace{2mm}
  
    \qrcode[hyperlink,height=0.75in]{https://www.menti.com/alqgaa5yhyd4}\\
    \vspace{2mm}
    \footnotesize  o ingresa a \href{https://www.menti.com/alqgaa5yhyd4}{menti.com}\\ con el código \textbf{5629 0857}
\end{frame}

\section{Objetivos}

\begin{frame}{Objetivos}
\begin{itemize}
    \item Brindar los conceptos necesarios para el diseño y manipulación de Bases de Datos relacionales.
    \item Explicar diferentes tipos de bases de datos y los criterios para escoger la más adecuada en el marco de un problema específico.
    \item Explicar los mecanismos que mejoran el desempeño de las Bases de Datos y las consecuencias de utilizar dichos mecanismos.
\end{itemize}
\end{frame}

\section{Contenido}

\begin{frame}{Contenido}
\begin{itemize}
    \item Modelo relacional
    \item SQL
    \item Diseño de Bases de Datos relacionales
\end{itemize}
\end{frame}

\section{Evaluación y Fechas}

\begin{frame}{Evaluación y Fechas}
    \begin{itemize}
        \item 20\% → Primer Parcial (Semana 6).
        \item 15\% → Segundo Parcial (Semana 12).
        \item 15\% → Examen Final (Semana 17).
        \item 20\% → Laboratorios (Semanales).
        \item 30\% → Proyecto Final (Segunda parte del semestre, 3 entregas).
    \end{itemize}
\end{frame}

\section{Reglas de Juego}

\begin{frame}{Reglas de Juego}
\begin{itemize}
    \item Trabajo Individual: Tareas, Parciales y Examen.
    \item Trabajo en Equipo: Talleres y Proyecto.
    \item Grupos de 2 personas para talleres y hasta 4 personas para el proyecto.
\end{itemize}
\end{frame}

\end{document}
